\chapter*{Introduction}
\addcontentsline{toc}{chapter}{Introduction}
\markboth{Introduction}{Introduction}
\label{chap:introduction}


 	\noindent\lettrine[lines=2]{\textbf{D}}\\ ans un monde marqué par l’essor des technologies numériques, les applications mobiles occupent une place centrale dans la diffusion des savoirs et l’accompagnement des apprentissages. Le Saint Coran, en tant que référence spirituelle et culturelle majeure, n’échappe pas à cette dynamique. Plusieurs solutions numériques existent déjà pour la récitation et l’étude du Coran, mais elles s’appuient essentiellement sur la lecture Ḥafṣ ‘an ‘Āṣim, largement répandue dans le monde musulman.\\

Or, dans le contexte maghrébin, la lecture Qālūn ‘an Nāfi‘ est la référence traditionnelle et constitue un patrimoine immatériel qu’il convient de préserver et de valoriser. L’idée du projet « Tarteel Maghreb » s’inscrit dans cette perspective : proposer une application mobile interactive dédiée à l’apprentissage de la récitation du Coran selon la lecture Qālūn, tout en intégrant des fonctionnalités d’intelligence artificielle pour l’analyse et l’amélioration de la prononciation.\\

Ce projet a été réalisé dans le cadre d’un stage d’été de deux mois, avec pour objectif de développer une application multiplateforme en React Native, appuyée par un back-end en Node.js et des services d’IA déployés via API. Il vise à offrir une expérience fluide et enrichissante aux utilisateurs, à la fois en ligne et hors-ligne, et à contribuer à la préservation de la tradition coranique maghrébine dans l’ère numérique.

    \vspace{0.5cm} % Espacement entre les paragraphes

    
	 
   \noindent 

